\begin{frame}{자주 묻는 질문 (1)}
  \textbf{Q. $k$는 몇으로?}
  \begin{itemize}
    \item $k=5$ 또는 $10$이 실전 기본값.
    \item $k$가 클수록($k=m$, \kb{LOOCV}) 각 fold 학습 데이터가 거의 전체 →
      편향 줄지만, fold를 $m$번 다시 학습 → 비용↑ · fold 간 분산↑.
    \item $k$가 작을수록($k=3$)은 빠르지만 fold 학습 데이터 작아 편향↑.
    \item $5\sim10$이 실용적 절충점.
  \end{itemize}
  \vskip0.4em
  \textbf{Q. 교차검증은 항상 해야 하나요?}
  \begin{itemize}
    \item 데이터가 아주 많으면(수백만 개) 한 번만 train/val로 나눠도 검증셋이
      충분히 커 안정적인 평가.
    \item 적을수록(수백$\sim$수천) 한 번의 분할이 우연히 ``쉬운/어려운''
      검증셋을 뽑을 위험이 커 → 교차검증의 가치↑.
  \end{itemize}
\end{frame}
